% \section{Conclusion}
% % \todo{Conclude after the cross-level synthesis and outlook are complete.}

% This survey examined recursive self-improvement through the scope of responsibility that AI systems assume for their own improvement. Taking the improvement loop as the unit of analysis, we organized the literature into five autonomy levels: executing prescribed updates, selecting improvement strategies, directing future learning experience, adapting persistently to deployment environments, and revising mechanisms that govern subsequent improvement. Across these levels, three questions anchor the analysis: where the loop closes, what changes are inherited, and which decisions remain externally controlled. We connected this framework to scientific discovery, embodied intelligence, software engineering, and industrial practice, as well as healthcare, while distinguishing demonstrated mechanisms from emerging research visions.

% The reviewed evidence establishes meaningful progress toward persistent and increasingly autonomous improvement, but remains uneven across mechanisms and application domains. Current systems can retain useful changes to parameters, memory, skills, and agent implementations; some also revise and reuse improvers, evaluators, or research policies. However, sustained and transferable gains in the capacity to improve remain insufficiently established. Greater autonomy must therefore be assessed separately from improvement quality and resource efficiency. In particular, inheriting a revised improvement mechanism establishes a recursive structure, while demonstrating its effectiveness requires evidence that it produces better subsequent improvements under comparable resources and independent evaluation.

% Progress toward genuine RSI depends on connecting useful experience acquisition, reliable validation, and persistent inheritance across successive rounds. Evaluation must establish whether gains survive changing tasks, preserve earlier capabilities, and justify the computation and human effort required to obtain them. As improvement mechanisms become adaptive, protected evaluation and explicit authority over objectives, resources, and deployment remain essential to interpreting and governing their effects. The central research objective is to build systems whose accumulated experience makes subsequent improvement more effective, efficient, and capable of discovering useful advances beyond prescribed strategies. Achieving this would turn today's bounded improvement loops into a dependable foundation for sustained AI development.

\section{Conclusion}

In this report, we present an autonomy-centered roadmap for recursive self-improvement, spanning five stages: improvement execution, strategy selection, experience acquisition, environmental adaptation, and recursive meta-improvement. Taking the improvement loop as the unit of analysis, we clarify what systems change, what successors inherit, and which decisions remain human-controlled. We connect this roadmap to scientific discovery, embodied intelligence, software engineering, and healthcare, showing how feedback availability, verification costs, and deployment constraints shape progress toward RSI in each domain. Industry practices and preliminary empirical findings further ground the roadmap in demonstrated capabilities and current limitations. Long-horizon evaluation is needed to establish whether these mechanisms deliver transferable gains under explicit resource constraints and human oversight. Ultimately, the promise of RSI lies in enabling each generation of AI to make future improvement more reliable, efficient, and conducive to novel discoveries.

%The reviewed evidence shows meaningful progress toward persistent, autonomous improvement, while sustained gains in the capacity to improve remain insufficiently established. Advancing toward genuine RSI requires reliable diagnosis, useful learning experience, trustworthy verification, and effective inheritance across successive rounds. 



%This survey examined recursive self-improvement through the scope of responsibility that AI systems assume for their own improvement. Taking the improvement loop as the unit of analysis, we organized the literature into five autonomy levels. Across these levels, three questions anchor the analysis: where the loop closes, what changes are inherited, and which decisions remain externally controlled.  The reviewed evidence establishes meaningful progress toward persistent and increasingly autonomous improvement. Current systems can retain useful changes to parameters, memory, skills, and agent implementations; some also revise and reuse improvers, evaluators, or research policies. However, sustained and transferable gains in the capacity to improve remain insufficiently established. We connected this framework to various scenarios, together with industrial practice, while distinguishing demonstrated mechanisms from emerging research visions.


%Greater autonomy must therefore be assessed separately from improvement quality and resource efficiency. In particular, inheriting a revised improvement mechanism establishes a recursive structure, while demonstrating its effectiveness requires evidence that it produces better subsequent improvements under comparable resources and independent evaluation.

%The central potential of RSI lies in making the production of further improvements more efficient. Accumulated experience can guide systems toward more informative learning opportunities, better-targeted interventions, and more effective validation, while inherited knowledge and strategies can reduce repeated search, diagnosis, and rework. Investments in data, computation, and human expertise can thereby benefit multiple improvement rounds, allowing useful discoveries to become reusable foundations for subsequent progress. This yields a concrete objective: achieving more durable, validated capability gains within a given total budget, or reducing the marginal cost of reaching a comparable capability target. Realizing this potential may still require substantial initial training and infrastructure; the efficiency advantage must emerge from how effectively these investments support later improvements.

%Progress toward genuine RSI should therefore be evaluated across successive rounds, accounting for the costs of experience acquisition, candidate search, validation, deployment, and human review. Gains must transfer to changing tasks, preserve earlier capabilities, and remain credible under protected evaluation and explicit authority over objectives, resources, and deployment. The research goal is to build systems that discover useful advances beyond prescribed strategies and retain the knowledge and mechanisms that make further advances easier to obtain. Ultimately, the promise of RSI is to turn accumulated experience into a progressively more productive capacity to improve, enabling sustained capability growth through reusable advances and lower marginal demands on computation and human effort.